
\section{Multiplication}

In this lesson, we will review multiplication. Multiplication is used when we \textit{repeatedly add}.
For example, we will find the number of blocks in this rectangle:
\\[0.1in]
\begin{tightcenter}
   \begin{tikzpicture}
      \draw[step=1.0cm,gray,thick]
      (0,0) grid (4,3);
   \end{tikzpicture}
   \captionof{figure}{a $4$ by $3$ grid}
\end{tightcenter}

We can see each row has $4$ unit squares. To find the total number of unit squares, we could do:
\[
   4+4+4=12
\]
Or, we use multiplication:
\[
   4 \times 3 = 12
\]
It's obvious multiplication is easier.\\*
The result of multiplication is called the \textit{product}. In the equation above, $12$ is the product of $4$ and $3$.\\*
This is a good opportunity to review these two important words:
\begin{itemize}
\item The result of addition is called the \textit{sum}. For example, in $1+2=3$, the number $3$ is the sum of $1$ and $2$.
\item The result of subtraction is called the \textit{difference}. For example, in $3-2=1$, the number $1$ is the difference of $3$ and $2$.
\end{itemize}
In later math courses, we regularly use the variable $x$, which causes confusion between $x$ and the multiplication symbol $\times$. To avoid such confusions, we use $\cdot$ to replace the multiplication symbol. We will re-write the last equation as:
\[
   4 \cdot 3 = 12
\]

It is critical to memorize the multiplication table. If you cannot quickly recall multiplication facts yet, please spend some time every day to memorize multiplication facts in \cref{tab:multiply}.

\begin{landscape} 
\thispagestyle{empty}
\begin{centering}
\large{\textbf{Multiplication Table}}\\
\end{centering}
\medskip

Here is how to use this table: On the first day, read "one times one is one", "one times two is two", all the way to "one times twelve is twelve."

On the next day, in addition to the row above, read from "two times one is two" all the way to "two times twelve is twenty-four."

Add in one row each day. Read as loud as possible. After two to three months, you will be able to quickly recall these multiplication facts. You cannot be successful in math courses without this ability.
\medskip

\footnotesize
\begin{tabular}{>{$}c<{$}|*{12}{>{$}c<{$}}}
    ~   &1&2&3&4&5&6&7&8&9&10&11&12   \\
    \hline\vrule height 12pt width 0pt

\\ 1&1\cdot 1=1&1\cdot 2=2&1\cdot 3=3&1\cdot 4=4&1\cdot 5=5&1\cdot 6=6&1\cdot 7=7&1\cdot 8=8&1\cdot 9=9&1\cdot 10=10&1\cdot 11=11&1\cdot 12=12 \\

\\ 2&2\cdot 1=2&2\cdot 2=4&2\cdot 3=6&2\cdot 4=8&2\cdot 5=10&2\cdot 6=12&2\cdot 7=14&2\cdot 8=16&2\cdot 9=18&2\cdot 10=20&2\cdot 11=22&2\cdot 12=24 \\

\\ 3&3\cdot 1=3&3\cdot 2=6&3\cdot 3=9&3\cdot 4=12&3\cdot 5=15&3\cdot 6=18&3\cdot 7=21&3\cdot 8=24&3\cdot 9=27&3\cdot 10=30&3\cdot 11=33&3\cdot 12=36 \\

\\ 4&4\cdot 1=4&4\cdot 2=8&4\cdot 3=12&4\cdot 4=16&4\cdot 5=20&4\cdot 6=24&4\cdot 7=28&4\cdot 8=32&4\cdot 9=36&4\cdot 10=40&4\cdot 11=44&4\cdot 12=48 \\

\\ 5&5\cdot 1=5&5\cdot 2=10&5\cdot 3=15&5\cdot 4=20&5\cdot 5=25&5\cdot 6=30&5\cdot 7=35&5\cdot 8=40&5\cdot 9=45&5\cdot 10=50&5\cdot 11=55&5\cdot 12=60 \\

\\ 6&6\cdot 1=6&6\cdot 2=12&6\cdot 3=18&6\cdot 4=24&6\cdot 5=30&6\cdot 6=36&6\cdot 7=42&6\cdot 8=48&6\cdot 9=54&6\cdot 10=60&6\cdot 11=66&6\cdot 12=72 \\

\\ 7&7\cdot 1=7&7\cdot 2=14&7\cdot 3=21&7\cdot 4=28&7\cdot 5=35&7\cdot 6=42&7\cdot 7=49&7\cdot 8=56&7\cdot 9=63&7\cdot 10=70&7\cdot 11=77&7\cdot 12=84 \\

\\ 8&8\cdot 1=8&8\cdot 2=16&8\cdot 3=24&8\cdot 4=32&8\cdot 5=40&8\cdot 6=48&8\cdot 7=56&8\cdot 8=64&8\cdot 9=72&8\cdot 10=80&8\cdot 11=88&8\cdot 12=96 \\

\\ 9&9\cdot 1=9&9\cdot 2=18&9\cdot 3=27&9\cdot 4=36&9\cdot 5=45&9\cdot 6=54&9\cdot 7=63&9\cdot 8=72&9\cdot 9=81&9\cdot 10=90&9\cdot 11=99&9\cdot 12=108 \\

\\ 10&10\cdot 1=10&10\cdot 2=20&10\cdot 3=30&10\cdot 4=40&10\cdot 5=50&10\cdot 6=60&10\cdot 7=70&10\cdot 8=80&10\cdot 9=90&10\cdot 10=100&10\cdot 11=110&10\cdot 12=120 \\

\\ 11&11\cdot 1=11&11\cdot 2=22&11\cdot 3=33&11\cdot 4=44&11\cdot 5=55&11\cdot 6=66&11\cdot 7=77&11\cdot 8=88&11\cdot 9=99&11\cdot 10=110&11\cdot 11=121&11\cdot 12=132 \\

\\ 12&12\cdot 1=12&12\cdot 2=24&12\cdot 3=36&12\cdot 4=48&12\cdot 5=60&12\cdot 6=72&12\cdot 7=84&12\cdot 8=96&12\cdot 9=108&12\cdot 10=120&12\cdot 11=132&12\cdot 12=144 \\

\end{tabular} 
\captionof{table}{Multiplication Table}
\label{tab:multiply}
\normalsize
\end{landscape}

Next, let's look at some examples where multiplication is used.
\thispagestyle{fancy}
\begin{myexample}
   Tim does $7$ math exercises every day. How many math exercises would Tim do in $9$ days?
\end{myexample}
\begin{solution}
	In $9$ days, Tim would do a total of $7 \cdot 9 = 63$ math problems.
\end{solution}

In this problem, we could have done $9 \cdot 7 = 63$. This is called the \textit{commutative property of multiplication}. Let's look at another example:
	\begin{itemize}
	\item $2 \cdot 3 \cdot 4 = 24$
	\item $2 \cdot 4 \cdot 3 = 24$
	\item $4 \cdot 3 \cdot 2 = 24$	
	\end{itemize}
In multiplication, we can change the order of numbers, and the product will not change.

Let's look at more multiplication examples.

\begin{myexample}
   Noel pays \$$50$ per month for a company to mow her lawn. How much money would Noel pay the company every year?
\end{myexample}
\begin{solution}
	There are 12 months in a year. If Noel pays \$$50$ per month, she would pay the company \$$50 \cdot 12 =$ \$$600$ every year. 
\end{solution}

\begin{myexample}
	Each crate has $10$ boxes of coke, and each box has $12$ cans. A shop ordered $50$ crates of coke. How many cans of coke did the shop order?
\end{myexample}
\begin{solution}
Since each crate has $10$ boxes of coke, $50$ crates have a total of $10 \cdot 50 = 500$ boxes.

Since each box has $12$ cans, $500$ boxes have a total of $12 \cdot 500 = 6000$ cans.

We could do this problem in one step:
	\[ 10 \cdot 50 \cdot 12 = 500 \cdot 12 = 6000 \]
	
\textbf{Conclusion:} The shop ordered $6,000$ cans of coke.
\end{solution}

\subsection{Summary}
Let's review important contents we learned in this lesson:

\begin{itemize}
\item We will stop using the multiplication symbol and use a dot instead. For example, $3 \times 4$ should be written as $3 \cdot 4$.
\item The result of addition is called the \textit{sum}. For example, in $1+2=3$, the number $3$ is the sum of $1$ and $2$.
\item The result of subtraction is called the \textit{difference}. For example, in $3-2=1$, the number $1$ is the difference of $3$ and $2$.
\item The result of multiplication is called the \textit{product}. For example, in $3 \cdot 4 = 12$, the number $12$ is the product of $3$ and $4$.
\item If the order of multiplication is changed, the product doesn't change. For example, $3 \cdot 4 = 4 \cdot 3$, and $2 \cdot 3 \cdot 4 = 3 \cdot 4 \cdot 2$.
\end{itemize}
